The baseline ADM analysis assumes an insurance coverage ratio of $\iota = 0.6$ (Section~\ref{sec: adm}). Because ADM does not publicly disclose facility-level property insurance terms, we report the sensitivity of key outputs to this assumption. Table~\ref{tab: insurance sensitivity} varies $\iota$ from 0.4 (conservative, lower coverage) to 0.8 (high coverage, typical for well-insured Fortune~500 operators).

\begin{table}[ht]
\centering
\small
\begin{tabular}{|c|c|c|c|c|c|}
\hline
$\iota$ & CapEx$_{\text{net}}$ (\$M) & $r_{\text{OpEx}}$ & OpEx (\$M) & $\ECL$ (\$M) & $\ICR_C$ \\
\hline
0.4 & 200.7 & 0.18 & 36.1 & 236.8 & 5.44 \\
0.5 & 167.2 & 0.18 & 30.1 & 197.4 & 5.45 \\
\textbf{0.6} & \textbf{133.8} & \textbf{0.18} & \textbf{24.1} & \textbf{157.9} & \textbf{5.46} \\
0.7 & 100.4 & 0.18 & 18.1 & 118.4 & 5.47 \\
0.8 &  66.9 & 0.18 & 12.0 &  78.9 & 5.48 \\
\hline
\end{tabular}
\caption{Sensitivity of ADM loss estimates to the insurance coverage ratio $\iota$. The bold row corresponds to the baseline assumption. $r_{\text{OpEx}}$ is insensitive to $\iota$ because it depends on $\DF_{\text{sector}}$, not on the net loss. $\ICR_C$ varies only modestly because ADM's EBIT (\$3.539\,B) is large relative to the annualised expected loss.}
\label{tab: insurance sensitivity}
\end{table}

The per-event $\ECL$ ranges from \$79\,M ($\iota = 0.8$) to \$237\,M ($\iota = 0.4$), a factor of $3\times$, while the climate-adjusted $\ICR_C$ remains within the A3/A- band across all scenarios. The credit risk conclusion -- that a 100-year flood event does not trigger a rating transition for ADM -- is robust to the insurance assumption. The \textit{magnitude} of the expected loss, however, is sensitive to $\iota$ and should be interpreted with the appropriate uncertainty in downstream applications.
